\section{Reusable domain templates}
\label{sec:domains}

The framework is useful only if domain experts can recognize where their
ordinary proof enters it.  This section gives conversion templates rather than
claiming that every deep producer theorem has already been proved.

\subsection{Finite algebra and linear algebra}

\paragraph{Objects.}
Coded commutative rings, modules, matrices, polynomials, finite presentations,
and bounded monomial tables.

\paragraph{Primitive-recursive core.}
Matrix arithmetic, determinants by Leibniz or fraction-free elimination,
characteristic polynomials, polynomial evaluation, truncated power series,
finite Gröbner/Macaulay calculations at a supplied degree, and verification of
factorizations or syzygies.

\paragraph{Typical contracts.}
A matrix inverse consumes a determinant/pivot margin or an exact adjugate
certificate.  A spectral projector consumes a gap separating the target
cluster.  Polynomial roots are returned as clusters unless a separation
certificate permits labels.

\paragraph{What becomes straightforward.}
Fixed-dimension identities and variable-dimension identities with bounded
folds are ideal trace-compiler targets.  Cayley--Hamilton, determinant
identities, finite module calculations, and exact certificate checking all
reduce to polynomial identities and bounded sums.

\subsection{Constructive analysis}

\paragraph{Objects.}
Regular rational names, interval or ball names, finite coefficient tables,
moduli of continuity, derivative bounds, and tail estimates.

\paragraph{Primitive-recursive core.}
Rational interval arithmetic, finite sums and products, Taylor polynomials,
quadrature rules, Euler--Maclaurin sums, bisection, and finite contour meshes.

\paragraph{Contracts.}
Every limiting theorem must expose a modulus or a finite error majorant.  A
partial operation such as reciprocal or logarithm consumes an apartness or
positivity margin.  An oscillatory or branch-sensitive construction returns a
relation until a gap resolves it.

\paragraph{Residual producer obligations.}
A classical existence proof of a convergence constant does not automatically
produce a primitive-recursive bound.  The framework can verify any proposed
constant and compose its consequences, but it cannot erase ineffectivity.

\subsection{Number theory and Diophantine geometry}

\paragraph{Objects.}
Integers, rationals, ideals in finite presentations, bounded-height points,
finite descent data, local solubility certificates, and exact recurrence
sequences.

\paragraph{Primitive-recursive core.}
Divisibility, gcd, modular arithmetic, continued fractions, finite lattice
search, bounded rational-point enumeration, verification of descent maps, and
finite prime sums.

\paragraph{Deep interface.}
A Diophantine finiteness theorem is split into:
\begin{enumerate}
\item a producer of an explicit height bound or complete list;
\item a bounded enumerator and verifier;
\item a proof that the list is complete below the bound.
\end{enumerate}
The last two are generic and iPRA-friendly.  The first contains the deep
mathematics.  For Mordell--Weil descent, an abstract finite Selmer group is not
enough: rational representatives on covering spaces, or a relational
certificate explaining their absence, must be supplied.

\subsection{Analytic number theory}

\paragraph{Objects.}
Finite prime tables, rational approximations to logarithms and exponentials,
truncated Dirichlet sums, zero-free-region certificates, and explicit error
budgets.

\paragraph{Primitive-recursive core.}
Sieve tables, divisor sums, partial summation, finite Euler products,
Möbius transforms, and verification of contour or kernel estimates.

\paragraph{Quantifier discipline.}
A useful theorem states the exact dependence of the cutoff on all numerical
parameters.  Phrases such as ``for sufficiently large \(x\), uniformly in
\(q\)'' must be expanded into a primitive-recursive bound with the quantifier
order fixed before compilation.

\paragraph{Residual obligations.}
Ineffective exceptional-zero arguments or unspecified constants remain
producer obligations.  The framework can still formalize conditional versions
and every finite verification step.

\subsection{Combinatorics}

Finite combinatorics is the most direct domain.  Graphs, set partitions,
finite posets, words, tableaux, and recurrences are bounded data.  Enumeration,
dynamic programming, involutions, sign-reversing pairings, and generating-
function coefficient identities compile to bounded traces.  Asymptotic claims
need explicit remainder estimates, but exact recurrences and congruences do
not.

For example, Bell-number recurrences and Touchard congruences are finite fold
theorems.  Dobi\'nski-type formulas require a truncation bound and a rational
tail estimate.  The interface cleanly separates the exact combinatorial core
from the analytic producer.

\subsection{Probability}

\paragraph{Objects.}
Finite rational probability laws, bounded random variables, rational moments,
finite product spaces, and explicit tail or quantization bounds.

\paragraph{Primitive-recursive core.}
Expectation is a finite weighted sum.  Independence is a table identity.
Convolution, characteristic-function samples, and finite couplings are
bounded calculations.

\paragraph{Approximation interface.}
A general distribution is represented by finite rational laws at precision
\(p\), together with a transport, tail, or anti-concentration modulus.  CDF
comparisons at discontinuities are relational or guarded.  A central-limit
statement becomes a finite estimate between rational laws plus explicit
truncation and smoothing errors.

\subsection{Topology}

\paragraph{Objects.}
Finite simplicial or cubical complexes, subdivision tables, chain complexes,
carrier clouds, finite covers, and quantitative compact metric presentations.

\paragraph{Primitive-recursive core.}
Boundary matrices, chain maps, parity, Sperner labels, simplicial subdivision,
finite homology certificates, and explicit prism homotopies.

\paragraph{Continuity interface.}
A map is observed through a modulus assigning a sufficiently fine source
level to each target level.  The finite map may be a cloud of target vertices.
A strong carrier condition makes every selection simplicial and any two
selections contiguous, so the induced homology map is independent of choice.
This resolves the noncanonical-selection issue without producing a global
continuous selector.

\subsection{Algebraic geometry}

Algebraic geometry can be split into exact finite algebra and genuinely
geometric producer theorems.

\paragraph{Affine presentations.}
An affine scheme of finite type is represented by a polynomial presentation
\(A=R[x_1,\ldots,x_n]/I\) together with bounded algebraic certificates.  A
morphism is a tuple of polynomials satisfying the target relations modulo
\(I\), witnessed by explicit combinations.  Equality of morphisms on a finite
presentation is another ideal-membership certificate.

\paragraph{Standard opens and gluing.}
A finite affine cover by standard opens \(D(f_i)\) carries an explicit cover
identity, for example \(\sum_i a_i f_i^N=1\) for a supplied common exponent
\(N\).  Local sections are fractions with synchronized
denominators.  Overlap compatibility is checked by cross multiplication after
a supplied localization exponent.  The finite affine-gluing theorem in
\cref{sec:casestudies} then constructs or verifies a global section.

\paragraph{Coherent sheaves.}
On a fixed finite affine cover, a finitely presented module is a matrix.  A
morphism is a matrix satisfying relation equations.  Kernels, cokernels, and
\v{C}ech complexes at a fixed degree are finite matrix calculations.  Exactness
is certified by explicit preimages or rank/minor data appropriate to the base
ring.

\paragraph{Points and dimensions.}
Bounded-height rational points are finite enumerations.  Dimension at a fixed
finite presentation can be certified through chains of prime ideals or finite
Gröbner data when an explicit degree bound is available.  General uniform
degree bounds and rational-point finiteness remain producer theorems, not
consequences of the interface alone.

\paragraph{Formal schemes and limits.}
An inverse system is represented only through finite truncations and explicit
transition maps.  A theorem about a requested truncation compiles directly.
A completed inverse-limit point requires a coherence/diameter theorem; it is
not obtained from levelwise nonemptiness alone.

\subsection{Category theory}

Finite categories, functors, natural transformations, adjunction data, and
finite diagrams are bounded tables.  Yoneda components, finite limits and
colimits, and finite Kan extensions reduce to table equations and bounded
search.  For categories of infinite mathematical structures, use categories
of finite presentations and contracted morphisms.  A categorical theorem is
then a reusable compiler for finite diagrams, not a quantification over an
uncoded universe of all objects.

\subsection{Differential and Riemannian geometry}

A charted object is represented by finitely many rational charts, overlap
margins, transition approximations, and derivative bounds.  Tensor identities
at rational sample points are finite algebra.  Curvature or geodesic
statements over a compact region require moduli and finite nets.  Coordinate
singularities are handled by relational chart clouds, while exact chart
selection requires a positive distance-to-boundary margin.

\subsection{Domain conversion table}

\begin{longtable}{@{}p{0.15\textwidth}p{0.23\textwidth}p{0.28\textwidth}p{0.22\textwidth}@{}}
\toprule
Domain & Finite presentation & Critical quantitative input & Safe output \\
\midrule
\endfirsthead
\toprule
Domain & Finite presentation & Critical quantitative input & Safe output \\
\midrule
\endhead
Algebra & coefficient/matrix tables & degree, pivot, or separation bound & exact certificate or cluster \\
Analysis & rational names and intervals & modulus, tail, apartness margin & regular name or guarded result \\
Diophantine geometry & bounded-height codes & height bound / covering representatives & complete finite list or conditional list \\
Analytic number theory & finite sums and kernels & effective cutoff and constants & explicit inequality \\
Probability & finite rational laws & tail and anti-concentration bounds & finite law approximation \\
Topology & finite complexes/carriers & subdivision and mesh modulus & chain/homology certificate \\
Algebraic geometry & finite presentations and affine covers & degree/localization/height bounds & finite algebra certificate or relational section \\
Category theory & finite categories/diagrams & finite bounds only & exact table theorem \\
Geometry & charts and rational nets & overlap and derivative margins & approximate invariant / guarded chart \\
\bottomrule
\end{longtable}
